\section{Discussion}
\paragraph{Why the loss collapses for discrete depth.}
When the block $f_\theta$ is applied $D$ times at inference with
$D\neq D_{\text{train}}$, the residual stream traverses a trajectory
unseen during training. Cross-entropy explodes catastrophically. The ODE
solver, by contrast, sees the same endpoint $h(1)$ regardless of solver
detail (to first order, ignoring discretisation error), so it is
inherently robust.

\paragraph{When does the latency advantage flip?}
The latency advantage of discrete depth in the very small regime
(<100k parameters) stems from Python-level overhead of \texttt{odeint}.
At the scale of the \texttt{gpu\_full.yaml} configuration
($d_{\text{model}}=256$, sequence length 64) the per-step compute
dominates, and the ODE solver's adaptive step control lets us pay only
for the precision actually required. This is the regime where the
authors of \cite{chen2018node} originally claimed wall-clock advantages.

\paragraph{Limitations.}
(i) Single synthetic task; (ii) tiny model; (iii) we use
backward-through-solver, not adjoint, so memory grows with NFE; (iv) no
analysis of training-time depth schedules (a richer continuation of this
work).
